% Lösungen Einheit 3 von 4 – Einheit 3 (zusammenbau v0.1)
\einheitenkopf{Einheit 3 von 4 · Einheit 3}

\erg{15}{a) mit Faktor ein Drittel – der Körper hat eine Spitze (Pyramide) \quad b) $G = 16$, $h = 6$, $V = \tfrac{1}{3} \cdot 16 \cdot 6 = 32$ \quad c) OS steht senkrecht auf der Grundfläche, ist also die Höhe; $G = \tfrac{1}{2} \cdot 5 \cdot 6 = 15$; $\tfrac{1}{3} \cdot 15 \cdot |OS| = 40$, $|OS| = 8$ \quad d) Diagonalen $|PR| = 7$, $|QT| = 6$, $G = \tfrac{1}{2} \cdot 7 \cdot 6 = 21$, $V = \tfrac{1}{3} \cdot 21 \cdot 5 = 35$ \quad e) $\overrightarrow{PP'} = \overrightarrow{QQ'} = \overrightarrow{RR'} = (0 | 0 | 6)$ steht senkrecht auf der Grundfläche in der $x$-$y$-Ebene; $G = \tfrac{1}{2} \cdot 4 \cdot 3 = 6$, $V = 6 \cdot 6 = 36$ \quad f) Grundkante 8, Höhe 3, Seitenhöhe $= \sqrt{4^2 + 3^2} = 5$; $O = 64 + 4 \cdot \tfrac{1}{2} \cdot 8 \cdot 5 = 144$ \quad g) $\overrightarrow{PQ} = \overrightarrow{TR}$ und $\overrightarrow{PQ} \cdot \overrightarrow{QR} = 0$: Rechteck; $|PQ| = 5 \ne |QR| = 10$: kein Quadrat; $G = 50$, $h = 6$, $V = \tfrac{1}{3} \cdot 50 \cdot 6 = 100$}
\erg{16}{a) Lena vergisst den Faktor ein Drittel der Pyramide. Richtig: $V = \tfrac{1}{3} \cdot 24 \cdot 5 = 40$ \quad b) Ein Würfel zerfällt in drei gleiche Pyramiden mit gemeinsamer Spitze in einer Ecke, jede mit einer Würfelfläche als Grundfläche und der Kante als Höhe; nach Cavalieri gilt der Anteil ein Drittel für jede Pyramide mit gleicher Grundfläche und Höhe. \quad c) $G = \tfrac{1}{2} \cdot 24 \cdot 20 = 240\,\mathrm{m}^2$, $h = 18\,\mathrm{m}$, $V = \tfrac{1}{3} \cdot 240 \cdot 18 = 1\,440\,\mathrm{m}^3$, Bedarf $= 14{,}4 \cdot 0{,}5 = 7{,}2\,\mathrm{kW}$ – 8\,kW reichen}
